\section{Two-dimensional acyclic resolutions: independent audit}
\label{sec:E-audit}

In this section a bound-$\mu$ resolution of a simply connected space $W$
means a group $G$, an actual two-dimensional CW model $X=K(G,1)$,
and a map $t:X\to W$ inducing integral homology isomorphisms, with
every homomorphism from $G$ to a group of cardinality at most $\mu$
trivial. The trivial group is allowed. All groups and cell sets are small.
We use $\mathfrak c=2^{\aleph_0}$.

\paragraph{Statement E.1 (verified input for the cardinal argument).}
Let $P$ be a group and $w_1,\ldots,w_m\in P*F(t_1,\ldots,t_n)$.
If their exponent-sum rows have rank $m$ over $\mathbb Q$, then
\[
 P\longrightarrow (P*F(t_1,\ldots,t_n))/
                   \langle\!\langle w_1,\ldots,w_m\rangle\!\rangle
\]
is injective. This is precisely Theorem~1 of the OpenAI preprint
\emph{Nonsingular systems of equations over arbitrary groups}, dated
5 October 2026 \cite[Theorem~1 and Remark~5.1]{EOpenAI2026}.
The complete sixteen-page source was obtained from the
official repository and read; it imposes no restriction on $P$.
The source is a preprint, and no claim of formal verification is made here.

Here are the proof checks relevant to using that theorem. For unitary
variables of sizes $k_1\geq\cdots\geq k_r$, rational independence of
the signed block multiplicity rows supplies nested nonzero minors.
The incidence space parametrizing a unitary matrix and an $l$-plane
fixed by it has codimension $l^2$. Its Poincar\'e dual has leading
exterior term a nonzero multiple of the product of the first $l$
primitive cohomology generators of the unitary group. Restriction to
$X\mapsto X\oplus(-I)$ verifies this coefficient by a single transverse
intersection; shorter exterior terms cannot contribute to the
top-dimensional pairing. The pairing for the product of the unitary
variable groups is the product of the indicated nested minors, hence
nonzero. This establishes the required simultaneous fixed-space bound.

For the operator step, work in the matrix algebra over the group von
Neumann algebra, with its faithful, unnormalized trace $\tau$.
Writing $h(e^{2\pi it})=t-\lfloor t\rfloor$ and
$\ell(U)=\tau(h(U))$, differentiate
$\tau(f(Ue^{2\pi is h(V)}))$ for smooth $f$ on the circle.
Cyclicity of the trace gives the derivative
$\tau(f'(Ue^{2\pi is h(V)})h(V))$; it is at most $\tau(h(V))$
when $f'\leq1$. Smoothing the sawtooth from the right, including at
the cut $1$, gives $\ell(UV)\leq\ell(U)+\ell(V)$ by bounded
convergence. For a weighted cyclic shift of length $a$, the phases of
the $a$ roots give $(a-1)/2$ plus the phase of the product of its weights.
These facts remain valid for uncountable coefficient groups: every
square-summable coefficient vector has countable support, and the trace
is still faithful and normal.

For a planar band diagram, independence of the exponent rows forces
equal positive and negative disk counts for each relator type. If there
are $2K$ disks, $b$ bands and $q$ boundary components, then
$2-q=2K-b$. The fixed-space bound and subadditivity give the upper
phase estimate $3b/2-K$. Assign a small negative phase to each
nonexceptional boundary and the compensating positive phase to the
exceptional boundary. The weighted-cycle formula tends to the same
constant plus the nonnegative phase of its coefficient word. Faithfulness
therefore makes that word the identity. Finally, a finite normal-closure
expression for a coefficient kernel element gives a punctured disk;
disjoint level arcs of the variable-circle maps supply planar bands
on which the coefficient map is constant. Removing components from
the inside out proves coefficient injectivity. This checks the source's
comparison from a kernel relation to its planar obstruction, rather
than assuming that an arbitrary pairing of letters is planar.

\paragraph{Statement E.2 (confirmation of INDLOCAL A.19).}
Let $X=K(G,1)$ be a connected CW complex of dimension at most two, and put
\[
 \nu=\max\{\aleph_0,d(H_1(G;\mathbb Z)),
                          \operatorname{rank}H_2(G;\mathbb Z)\}.
\]
For each $g\ne1$ there is a homomorphism $G\to Q$ with
$|Q|\leq\nu^{\aleph_0}$ that does not kill $g$.

\emph{Proof.}
Start with a connected subcomplex of at most $\nu$ cells containing
a loop for $g$, cycles generating $H_1(X)$, and cycles giving a basis
of $H_2(X)$. This last group is free because it is a subgroup of
$C_2(X;\mathbb Z)$. At each of countably many stages, add finite
supports of nullhomotopies for every element of the kernel of the
current fundamental-group map to $G$, and finite bounding chains for
every element of the kernel of its first-homology map to $H_1(X)$.
There are at most $\nu$ choices at each stage. The union $Y$ has at
most $\nu$ cells, and the actual inclusion $Y\subseteq X$ induces
an injection on fundamental groups and an integral homology
isomorphism. Injectivity on $H_2$ requires no choice: there are no
three-boundaries in either complex.

Put $H=\pi_1Y$. Its embedding in $G$ identifies the universal cover
of $Y$ with a component of the inverse image of $Y$ in the
contractible universal cover of $X$. This component has no second
homology, since its cellular two-cycles inject into those of the
contractible two-complex. It is simply connected and
two-dimensional, hence contractible by Hurewicz and Whitehead.
Thus $Y$ itself is a classifying model, and its displayed subgroup
$H$ contains the specified nonidentity $g$.

Collapse a maximal tree of $Y^{(1)}$ and an extension of it to a
maximal tree of $X^{(1)}$. The resulting pair has one vertex and
the relative cellular boundary is an actual isomorphism
\[
 \delta:\mathbb Z^{(R)}\xrightarrow{\cong}\mathbb Z^{(S)},
\]
where $R,S$ are the relative two- and one-cells. For each $r\in R$
choose its finite attaching word over $H*F(S)$ and let $W(r)$ contain
\emph{all} its occurring variables, including those of total exponent
zero. For each $s\in S$, let $I(s)$ be the finite support of
$\delta^{-1}(s)$. A pair of subsets is closed when it is closed under
$r\mapsto W(r)$ and $s\mapsto I(s)$. Closing one additional cell
takes at most countably many cells. By well-ordering the cells this
gives a continuous filtration of actual subcomplexes from $Y$ to $X$
in which each successor adds countably many relative cells. On each
successor quotient $\delta$ remains an isomorphism. Van Kampen
therefore gives presentations whose exponent rows in the newly
added variables are independent, every finite subfamily included.

There is a group $E$ of cardinality at most $\nu$ containing $H$
and solving every finite nonsingular system over itself: perform
countably many passes through all finite systems over the group
present at the beginning of a pass, adjoining universal solutions.
Statement E.1 keeps all coefficient maps injective. Only $\nu$
many finite strings and generators occur in a pass.

Use the single reduced power
\[
 U=E^{\mathbb N}/E^{(\mathbb N)},\qquad
 E^{(\mathbb N)}=\{(e_i):e_i=1\text{ for all but finitely many }i\}.
\]
The denominator is normal; constant sequences embed $E$, and
$|U|\leq\nu^{\aleph_0}$. Any countable nonsingular system over $U$
is solvable in $U$: fix representatives of its coefficients, solve
the first $k$ equations in $E$ at coordinate $k$, and assemble the
solutions coordinatewise. Each equation holds at every sufficiently
large coordinate. Only the finitely many variables occurring in the
first $k$ equations have to be solved there, so E.1 applies exactly.

Now extend the displayed embedding $H\hookrightarrow U$ through
the filtration of the relative presentation. At each successor,
apply the current homomorphism to the coefficients and solve the
resulting countable system in $U$. Assigning the new variables to
these solutions gives a homomorphism on the next presentation with
\emph{exactly} the previous restriction. At a limit use the group
colimit universal property. The resulting homomorphism $G\to U$
is injective on $H$, so it detects $g$. Taking its image as $Q$
proves the assertion. No unproved injection of an intermediate
presentation group is used in this recursion.\hfill$\square$

\paragraph{Statement E.3 (confirmation of INDLOCAL A.20).}
Let $W$ be simply connected, two-dimensional, and
$H_2(W;\mathbb Z)\cong\mathbb Z^{(\kappa)}$ with $\kappa>0$.
There is no bound-$\mu$ two-dimensional resolution when
\[
 \mu\geq \max(\aleph_0,\kappa)^{\aleph_0}.
\]
Indeed any such group has $H_1=0$ and $H_2\cong\mathbb Z^{(\kappa)}$,
so is nontrivial. Statement E.2 detects any nonidentity element in
a quotient of the forbidden size. For $W=S^2$ the excluded range
is exactly $\mu\geq\mathfrak c$ from this argument; the argument
does not exclude $\aleph_0\leq\mu<\mathfrak c$.

\paragraph{Statement E.4 (an explicit finite-bound resolution).}
Every simply connected two-dimensional CW complex has a
two-dimensional resolution satisfying every finite bound at once.
For a target with finite rank $k>0$ in degree two, such a resolution
has a finite classifying model. A target of infinite rank $\kappa$
has one with at most $\kappa$ cells.

\emph{Proof.}
We give an explicit Higman-block proof of this case of the
finite-bound construction in \cite[Sections~8--10]{EBaumslagDyerHeller1980}.
Let
\[
 H=\langle a,b,c,d\mid a^{-1}ba=b^2,
 b^{-1}cb=c^2,\ c^{-1}dc=d^2,\ d^{-1}ad=a^2\rangle.
\]
Its presentation complex is aspherical. To check this, first
amalgamate the standard aspherical Baumslag--Solitar complexes on
$(a,b)$ and $(b,c)$ along the circle for $b$. The subgroup
$\langle a,c\rangle$ is free of rank two: nonzero powers of
$a$ and $c$ lie outside the amalgamated cyclic subgroup, so their
alternating words are reduced. Do the same on $(c,d,a)$, then
amalgamate the two resulting spaces along that free rank-two
subgroup. Injective amalgam normal forms identify the universal
cover as a tree of contractible spaces. The displayed four-relator
presentation is the associated two-dimensional model. Its ordinary
boundary $C_2\to C_1$ is a signed permutation matrix, so it is
acyclic, and each named generator has infinite order.

Every homomorphism $H\to F$ to a finite group is trivial. If any
generator image were nonidentity, choose the least prime $p$
dividing one of their orders, say the order of $x$. The preceding
image $y$ conjugates $x$ to $x^2$. Thus $x$ has odd order, and
the multiplicative order of $2$ modulo $p$ divides the order of
$y$. It is greater than one and divides $p-1$, contradicting
minimality of $p$.

For a cardinal $\kappa>0$, take one central copy of $H$ and
$\kappa$ leaf copies of $H$, amalgamating the infinite cyclic
subgroups generated by $a$ along the edges of a star. Denote
the resulting group by $G_\kappa$. Its graph-of-spaces model
uses only two-dimensional vertex spaces and cylinders on circles,
and is aspherical by the same normal-form argument. The
Mayer--Vietoris sequence for this tree gives
\[
 H_1(G_\kappa;\mathbb Z)=0,\qquad
 H_2(G_\kappa;\mathbb Z)\cong\mathbb Z^{(\kappa)},\qquad
 H_i(G_\kappa;\mathbb Z)=0\quad(i>2).
\]
Every finite image kills every vertex group, and hence all of
$G_\kappa$. The finite boundary matrices for the first eight stars
are independently checked in \texttt{code/E\_higman\_homology.py}.
This calculation is a homology check; asphericity is established
by the preceding geometric argument.

There is an actual homology equivalence to
$W_\kappa=\bigvee_\kappa S^2$. Collapse a maximal tree to make
the source one-vertex. Since $H_1=0$, its boundary
$C_2\to C_1$ is onto and splits over $\mathbb Z$. The splitting
gives a projection $C_2\to H_2$. Collapse the one-skeleton and
map each resulting two-sphere by the finite pinch-and-degree map
specified by its projected basis vector. This defines a continuous
cellular map to $W_\kappa$ and induces the chosen identity on
$H_2$. It is a homology isomorphism in every degree.
Finally every simply connected two-dimensional $W$ is homotopy
equivalent to $W_\kappa$: choose spherical representatives of a
basis of $H_2(W)$ using Hurewicz, form their wedge map, and apply
the simply connected homology Whitehead theorem. If $\kappa=0$,
$W$ is contractible and the trivial group suffices.\hfill$\square$

\paragraph{Statement E.5 (confirmation of INDLOCAL A.13).}
For every infinite cardinal $\lambda$ there is a simple group
$S_\lambda$ of size $\lambda$ with an actual two-dimensional
classifying space and a homology equivalence
$BS_\lambda\to\bigvee_\lambda S^2$. Consequently it is a
bound-$\mu$ resolution for every $\mu<\lambda$.

\emph{Proof.}
Start with $F_\lambda$. At each pass process every ordered pair
$g\ne1,h$ in the group at the beginning of the pass by adjoining
a free letter $t$ and an HNN letter $s$ satisfying
\[
 s^{-1}ws=wh,\qquad
 w=gtg^{-1}t^{-1}gt^2g^{-1}t^{-2}.
\]
In the free product with $\langle t\rangle$, $w$ is cyclically
reduced of syllable length eight. Cyclic reduction of $wh$
has length eight if $hg\ne1$ and six otherwise. Hence both
associated cyclic subgroups are infinite, and HNN normal form
embeds the old group. The relation puts $h$ in the normal
closure of $g$. Repeat the passes countably. Every pair of
elements occurs at a common pass, so the union is simple.

Perform the attachments on classifying spaces. A free letter adds
a circle; each HNN extension adds a circle and a two-cell.
These actual inclusions are aspherical graph-of-spaces extensions,
and their union is aspherical because every sphere map has finite
cell support. The group contains $F_\lambda$ and has at most
$\lambda$ generators, giving cardinality $\lambda$.

In the first pass there are $\lambda$ pairs $(g,1)$. Each of
their distinct two-cells has zero ordinary boundary. They remain
independent in $H_2$ because there are no three-cells. Simplicity
gives $H_1=0$; $H_2$ is free of rank exactly $\lambda$, and all
higher homology vanishes. The split-boundary construction in E.4
now gives the claimed actual map to the wedge. A nontrivial map
from a simple group is injective, which excludes targets of
size less than $\lambda$.\hfill$\square$

\paragraph{Audit E.6 (INDLOCAL A.10 and A.16).}
The general A.10 assertion uses the basic groupoid realization
of Baumslag--Dyer--Heller, Sections 8--10
\cite{EBaumslagDyerHeller1980}. Their Lemma 10.1 is
exactly the reduced-word injectivity needed for its boundary
inclusions, Theorem 9.1 converts the groupoid pushout to the
homotopy pushout, Lemma 10.2 gives the acyclic simplex pieces,
and Theorem 10.4 preserves the dimension for finite simplices.
After killing the inserted Higman factors, each pointed
simplex groupoid is indiscrete. All pointed intersections
contain the common base object; thus gluing those quotients
creates precisely the fundamental groupoid of the original
complex, with no extra loops. The kernel is consequently
normally generated by the inserted Higman groups, and all
finite-group maps kill it. Inductively attaching mapping
cylinders from the already constructed boundary models to
the finite simplex models preserves the required dimension
and cell bound. The local-coefficient comparison is obtained
simplex by simplex and then by cellular unions. These are
the steps needed to supplement the cited dimension theorem;
none is an assertion that arbitrary bar models preserve
dimension. E.4 provides a separate direct verification of
the simply connected two-dimensional cases needed here.

For A.16, the geometric input in A.14 also checks
\cite[Section~6]{EBartholdiMikhailov2026}. Let $L$
be a finite flag triangulation of a positive-genus surface
and form its real moment-angle cube complex $C(L)$. Remove
one open corner to get $M$. Its cut link is that surface.
The dual graph of the three-cells is connected; collapsing
along a spanning tree rooted at a cut face removes every
three-cell. The oriented cube sum maps to a generator in
$H_3(C(L),M)$, so the cut surface maps to zero in $H_2(M)$.
The radial retraction in the CAT(0) universal cover of
$C(L)$ proves injectivity of the old group when $M$ is
glued, through a surface cylinder, to an aspherical space.
After collapsing the lifted old spaces, the universal cover
has flag vertex links (covers of $L$ at the new vertices),
so the quotient is CAT(0). The relative homology sequence
and contractibility of the old lifted pieces make the
uncollapsed universal cover contractible too. The argument
does not require the old aspherical space to have dimension
two, which verifies the extension of the source lemma used
in A.16.

The algebraic steps of A.16 are likewise valid. In the
pullback cover over $\widetilde W$, let
$R=\mathbb Z[\pi_1W]$ and keep the initial copy
$B=H_2(\widetilde W)$. The simple-kernel extension adds
only relative cells in degrees one and two. Its surjective
relative boundary has projective kernel $Z_2$. If the old
$H_1$ has projective dimension at most one, then
$\ker(Z_2\to H_1)$ is projective, by the long exact
$\operatorname{Ext}$ sequence. Hence the new kernel of
$H_2\to B$ is projective. The deliberately included
$\lambda$ commutator relators with $h=1$ split off an
actual free $R^{(\lambda)}$ summand: their coordinates
occur in neither $\partial_2$ nor the old
$\operatorname{im}\partial_3$. The countable direct-sum
swindle makes the whole projective kernel free.

Attach a surface block for each of its basis vectors.
The map of that surface into $\widetilde W$ is nullhomotopic,
because it has zero $H_2$ class and the target is simply
connected. Thus the map to $W$ extends literally. Equivariant
Mayer--Vietoris kills precisely the chosen free kernel,
keeps homology above degree two unchanged, and leaves
$H_1$ of projective dimension at most one and the new
kernel of $H_2\to B$ free. Countable alternation gives
the desired simple kernel and homology equivalence on
pullback covers. The mapping cone of the cellular chain
comparison is an acyclic bounded-below complex of free
$R$-modules, so is contractible; tensoring proves the
comparison for every local system. The dimension bound
is $\max(\dim W,3)$ because the surface cylinders are
three-dimensional. All these steps confirm A.16 and
also identify why it does not answer E2.

\paragraph{Audit E.7 (INDLOCAL A.27--A.31).}
Here $[a,b]=aba^{-1}b^{-1}$. The following records the
individual checks, including the comparisons that preserve
two-dimensional classifying spaces.

For A.27's mixed-identity-free case, the commutator maps
$t\mapsto[a,t]b$ carry a word outside the coefficient
free factor to another word outside it, by malnormality
of a free factor in a free product. Their finite
composition minus the prescribed root value is therefore
a nonidentity mixed word and has a nonidentity evaluation.
This verifies the recursive coefficient choice in
(A.27.4). At a solution whose root equals its $k$th
enumerated value, the $k$th equation becomes constant
and backward substitution contradicts exactly that choice.
For a torsion-free two-conjugacy-class group,
$t\mapsto[a,t]b$ has image $E-\{ab\}$ with infinite
fibres; deleting finitely many inputs preserves its
image. This independently justifies the same recursive
choice in that case. For a uniformly simple group,
the word $W_a=\prod_{i=1}^N v_i[a,u_i]v_i^{-1}$ is
onto whenever $a\ne1$. Indeed a nontrivial commutator
$c=[a,u]$ normally expresses every element in at most
$N$ conjugates, and
$c^{-1}=u[a,u^{-1}]u^{-1}$. The finite tree
substitutions in A.27 are therefore onto, and the
chosen coefficients exclude each candidate root value.

Every displayed relator has exponent one in its own
variable and zero in all others. After lifting its
coefficients from the specified subgroup, finite-prefix
elimination of nonroot variables leaves one unimodular
root relator over the old torsion-free group free-product
with the remaining free letters. Forester--Rourke's
one-pair theorem \cite[Main Theorem]{EForesterRourkeDiagrams}
applies with exactly those hypotheses (an exponent-$\pm1$
shape is amenable):
the coefficient group embeds, and the inclusion-induced
extension-of-scalars map on $\pi_2$ is an isomorphism.
The prefix is therefore aspherical. Compactness of sphere
maps makes the full CW union aspherical; compactness
of nullhomotopies verifies injection of the old group
without needing injections between every two prefixes.
The identity relative boundary gives the integral
homology isomorphism. Killing the coefficient group
first kills the root and then all variables, proving
normal generation. A map extending the chosen subgroup
epimorphism would solve the forbidden system. This
checks the subgroup version required by A.28.

For A.28, a group of size at most $\mathfrak c$ has at
most $\mathfrak c^{\aleph_0}=\mathfrak c$ countable
subgroups. There are at most $\mathfrak c$ countable
group structures and maps, so a pass has at most
$\mathfrak c$ blocks. The $\omega_1$ outer passes
retain that bound. Every map onto a countable group
is onto on a countable subgroup lying in one earlier
pass, since $\operatorname{cf}(\omega_1)>\aleph_0$.
Its listed block contradicts extension of the final
map. This is the precise diagonalization; it never
assumes that a subgroup map extends to the whole
current group. The extra one-variable unimodular
adjunctions preserve asphericity and homology by the
same one-pair theorem, and every finite coefficient
list is eventually processed. Their solvability
passes to quotients simply by lifting coefficients.
Normal generation is transitive through the
attachments and unions.

The Higman image exclusions used in A.28 are valid.
The least-prime argument in E.4 works if the four
generator images merely have finite order, so kills
maps into any torsion group. In a linear image,
each generator is conjugate to its square, forcing
its eigenvalues to be roots of unity. In positive
characteristic it then has finite order. In
characteristic zero, after taking unipotent powers,
the identity
$\operatorname{Ad}(y^{-m})\log(x^L)=2^m\log(x^L)$
forces the logarithm to vanish: conjugation by a
unipotent matrix is unipotent. Thus every generator
again has finite order. A locally linear target
contains the finitely generated image in a linear
subgroup, so it too is killed. Normal generation
transfers these exclusions to the final group.
Mapping each new loop constantly to $S^2$ and
filling each attaching loop extends the original
map literally; the homology isomorphism follows
from the inclusion comparison.

For A.29, nested commutators combine a finite
equational cover of $E-\{1\}$ into a zero-exponent
word $v$ nonidentity at $1$ and identity elsewhere.
Then $v(x)x=1$ is an unsolvable unimodular equation.
Its finite-avoidance consequence supplies the
recursive weighted letters of the proposed metric.
At each stage there are finitely many constraints,
because their conjugators lie in finite balls of
the auxiliary proper length and their word lengths
are bounded. The stage invariant gives
$L(f_j)\geq2^{-j}>0$, while the chosen nonidentity
$a_i$ satisfy $L(a_i)\leq2^{-i}$. Conjugation
changes costs by at most the displayed finite
constant $D_g$, which proves continuity of the
group operations. Thus the claimed nondiscrete
Hausdorff metrizable topology follows. No
completeness or conjugation invariance is obtained.

For A.30, the conjugation-invariant metric does
give the uniform estimate
$\ell(v[b,u]v^{-1})\leq2\ell(b)$. A word with
$L$ variable occurrences is $L$-Lipschitz in the
maximum product metric. Since those conjugated
commutators form a symmetric normal generating
set, a finite path between two different word
values has a first changing step. Choosing $b$
small excludes the prescribed value throughout
the one-step image while retaining two distinct
values. Iteration on the binary tree gives nested
nonempty images excluding every group element.
For each finite prefix, eliminating parents
before children gives an actual homotopy
equivalence relative to the old space with its
wedge of frontier circles. This verifies its
asphericity and coefficient injection independently
of any simultaneous-equation asphericity conjecture.

A.31 replaces the metric estimate by the exact
word-norm bound
\[
 \rho_b\bigl(R(\mathbf s\mathbf u)R(\mathbf u)^{-1}\bigr)
 \leq2L,\qquad
 s_i\in\{v[b,u]v^{-1}\}.
\]
Replacing occurrences one at a time and multiplying
the consecutive differences proves the bound,
including inverse occurrences. Its hypothesis
allows $b$ for which the excluded difference
has greater norm, so the same binary-tree
construction works. If that hypothesis fails,
some $d\ne1$ has $\sup_b\rho_b(d)=M<\infty$.
Expressing any $g$ by $k$ conjugates of $d^{\pm1}$
gives $N_E(g)=\sup_b\rho_b(g)\leq kM$.
Moreover, for every fixed $d\ne1$,
\[
 \rho_d(g)\leq N_E(g)\leq N_E(d)\rho_d(g).
\]
Hence boundedness of any one $\rho_d$ would give
uniform simplicity. For the surviving quotients
all these norms are unbounded; pointwise finiteness
of $N_E$ is nevertheless consistent with this
calculation. A positive invariant length satisfies
$\ell(d)\leq N_E(d)\ell(b)$ for every $b\ne1$,
which proves the claimed discreteness. These
checks confirm the scope of A.31; they do not
eliminate its residual class.

\paragraph{Statement E.8 (the ranges established for E2--E4).}
For $\kappa=0$ every bound is possible. For finite
$\kappa>0$, E.3 and E.4 establish existence for
all finite bounds and nonexistence for
$\mu\geq\mathfrak c$. For infinite $\kappa$,
E.3 and E.5 establish existence for $\mu<\kappa$
and nonexistence for $\mu\geq\kappa^{\aleph_0}$.
In particular if $\kappa^{\aleph_0}=\kappa$,
the exact answer is $\mu<\kappa$.
The outstanding intervals in this work are
\[
 \aleph_0\leq\mu<\mathfrak c
 \quad(0<\kappa<\aleph_0),\qquad
 \kappa\leq\mu<\kappa^{\aleph_0}
 \quad(\kappa\text{ infinite}).
\]
No claim that these intervals are independent of
ZFC is made. Under CH the first interval contains
only $\aleph_0$; this arithmetic observation does
not decide its existence question. The smallest
unresolved case remains $W=S^2$, $\mu=\aleph_0$.

\paragraph{Statement E.9 (limitation of two-dimensional group tests).}
For any cardinal bound, groups with two-dimensional
classifying spaces cannot replace all the group
tests of INDLOCAL B.7 for arbitrary homotopy types.
Indeed, for every actual two-dimensional CW
complex $X$ and every $2$-connected space $Y$,
obstruction theory (or cellular approximation
for $Y=S^3$) gives $[X,Y]=\{*\}$.
In particular $[BG,S^3]$ is a singleton for every
such group, although $S^3$ is not a $1$-type.
Thus E2--E4 can address only the $S^2$ or
$\pi_2$ portion of a bounded-test refinement.
If the test category is the one-object group
$G$, its number of arrows is $|G|$. A nontrivial
bound-$\mu$ resolution usable among tests with
fewer than $\kappa$ arrows must satisfy
$\mu<|G|<\kappa$: the first inequality follows
from the identity homomorphism of $G$. The
missing existence assertion in E2 cannot be
inferred from B.7 or from this necessary size
condition.

\paragraph{Statement E.10 (a discarded proposed construction is impossible).}
There is no nontrivial simple locally free group.
More generally, a nontrivial locally free group
with a finite normal generating set cannot be
perfect. To prove this, choose nonidentity normal
generators $x_1,\ldots,x_r$. If the group were
perfect, write them as finite products of
commutators, then write the finitely many entries
of those commutators as products of conjugates
of the $x_i^{\pm1}$. All the finitely many
elements and conjugators lie in one finitely
generated free subgroup $F$. The Magnus
embedding of $F$ into noncommutative integral
power series assigns every nonidentity element
a least positive degree. Let $d$ be the least
such degree among the $x_i$. Conjugation and
products preserve degree at least $d$, while
the commutator of two such series has degree
at least $2d$. The chosen commutator expressions
therefore give degree at least $2d$ for every
$x_i$, contradicting minimality of $d$.
A nonabelian simple group is perfect and is
normally generated by any nonidentity element,
so the first assertion follows.

\paragraph{Statement E.11 (a restriction on cyclic HNN compression).}
Let $X$ be an actual two-dimensional $K(G,1)$ and
let $G\twoheadrightarrow\mathbb Z$ have kernel $N$.
Put $R=\mathbb Q[t,t^{-1}]$, with $t$ acting by a
generator of the deck group of the cyclic cover
$X_N\to X$. Then $H_2(N;\mathbb Q)$ is a free
$R$-module. If its $R$-rank is $\beta$, there is an
exact sequence
\[
 0\longrightarrow\mathbb Q^{(\beta)}
 \longrightarrow H_2(G;\mathbb Q)
 \longrightarrow
 \ker\bigl(t-1:H_1(N;\mathbb Q)\to H_1(N;\mathbb Q)\bigr)
 \longrightarrow0.
\]
Consequently $\beta\leq\dim_{\mathbb Q}H_2(G;\mathbb Q)$;
if $N$ is perfect, equality holds. Moreover,
\[
 \dim_{\mathbb Q}H_2(N;\mathbb Q)=
 \begin{cases}
 0,&\beta=0,\\
 \max(\aleph_0,\beta),&\beta>0.
 \end{cases}
\]
In particular, a perfect group with uncountable
second Betti number cannot be the cyclic normal
kernel of a two-dimensional group with finite
second Betti number.

For the proof, cellular chains of $X_N$ are free
$R$-modules, including when the sets of cells
are infinite. Since $X_N$ is two-dimensional,
$H_2(X_N;\mathbb Q)=\ker d_2$ is a submodule
of the free module $C_2(X_N;\mathbb Q)$.
The Laurent polynomial ring $R$ is a principal
ideal domain, so this submodule is free, with
some basis of cardinality $\beta$. The short
exact sequence of cellular chain complexes
\[
 0\longrightarrow C_*(X_N;\mathbb Q)
 \xrightarrow{\ t-1\ }C_*(X_N;\mathbb Q)
 \longrightarrow C_*(X;\mathbb Q)\longrightarrow0
\]
gives the displayed exact sequence in homology,
because the coinvariants of a free $R$-module
$R^{(\beta)}$ are $\mathbb Q^{(\beta)}$.
The underlying rational vector space of $R$ has
countably infinite dimension, proving the rank
formula. Finally $X_N$ is a $K(N,1)$, so these
are the group homology groups claimed.

Applied to $G=S\rtimes_f\mathbb Z$, this rules
out reducing an uncountable $H_2(S;\mathbb Q)$
to finite rank merely by choosing an automorphism
$f$ with small coinvariants while retaining an
actual two-dimensional model. For a proper
ascending HNN extension the cyclic kernel is
instead $\operatorname*{colim}(S\xrightarrow fS
\xrightarrow f\cdots)$; the same statement applies
to that kernel. Killing its homology first does
not produce a two-dimensional model for the
extension. If $S$ was simple, that direct limit
is simple, so an acyclic two-dimensional cyclic
kernel would already supply the difficult
uncountable acyclic simple group, rather than
deducing one from a homology calculation.

\paragraph{Statement E.12 (ordinary countable equation compactness fails).}
Let $E$ be a countably infinite group such that, for every $d\in E$,
there is a finite integer $n(d)$ for which $d$ is a product of at most
$n(d)$ conjugates of $a^{\pm1}$ for every $a\ne1$ in $E$.
Then there is a countable system of equations with coefficients in $E$
whose every finite subsystem has a solution in $E$, but which has no
simultaneous solution in $E$.
In particular the pointwise finite alternative for $N_E$ in
INDLOCAL A.31 cannot occur in a countably equation-compact group.

\emph{Proof.}
Fix a nonprincipal ultrafilter $\mathcal U$ on $\mathbb N$ and put
$P=E^{\mathbb N}/\mathcal U$, identifying $E$ with its diagonal copy.
For any $b\ne1$ in $P$ and any $d\in E$, the bound $n(d)$ implies
$d\in\langle\!\langle b\rangle\!\rangle_P$.
Indeed, choose representatives $b_i$. On a set in $\mathcal U$,
$b_i\ne1$, and there are expressions for $d$ as products of at most
$n(d)$ conjugates of $b_i^{\pm1}$.
There are only finitely many possible lengths and sign patterns;
one pattern occurs on a set in $\mathcal U$.
The conjugating elements in that pattern, chosen coordinatewise and
defined arbitrarily off this set, give the required expression in $P$.
Thus every nontrivial normal subgroup of $P$ contains the diagonal $E$.

Choose $x\in P\setminus E$, for instance the class of a sequence
of distinct elements of $E$. Construct countable subgroups
$P_0\le P_1\le\cdots$ of $P$, starting with $P_0=\langle E,x\rangle$.
At step $j+1$, for each $b\in P_j\setminus\{1\}$ and $d\in E$,
adjoin the finitely many conjugators in one expression witnessing
$d\in\langle\!\langle b\rangle\!\rangle_P$.
There are only countably many pairs, so
$P'=\bigcup_jP_j$ is countable and contains $E$ properly.
Every nontrivial normal subgroup of $P'$ contains $E$.

Use one variable for each element of $P'\setminus E$; for an element
of $E$ use that element itself as a coefficient.
Impose all the multiplication equations of $P'$.
This is a countable system. Every finite subsystem has a solution
in $P$, namely the displayed elements of $P'$, and consequently in
$E$: for the finitely many equations, take representatives of these
elements and choose one coordinate in the intersection of the
finitely many sets in $\mathcal U$ on which the equations hold.
If the entire system had a solution in $E$, its values would define
a homomorphism $r:P'\to E$ restricting to the identity on $E$.
Its kernel is nontrivial since
$x r(x)^{-1}\ne1$, whereas every nontrivial normal subgroup of $P'$
contains $E$. This contradicts $r|_E=\mathrm{id}_E$.
\hfill$\square$

The exponent-sum rows of this multiplication-table system are not
asserted independent. The argument supplies neither a nonsingular
system nor an actual two-dimensional relative homology equivalence.
Adding fresh variables $y_i=r_i$ to make arbitrary equations
formally independent removes their restrictions; imposing $y_i=1$
restores the dependency. Thus E.12 gives a precise compactness
obstruction, but does not decide E2.

\paragraph{Statement E.13 (the independent geometric obstruction).}
An invertible integral exponent matrix alone does not make a
finite relative presentation aspherical, even over the trivial
torsion-free coefficient group. For example, the presentation
\[
 P=\langle a,b\mid a^5=b^3,\ b^3=(ba)^2\rangle
\]
defines the binary icosahedral group of order $120$
\cite[Example~3.2]{EMikhailovOrr2013}.
Its exponent matrix is
\[
 \begin{pmatrix}5&-3\\-2&1\end{pmatrix},
 \qquad \det=-1.
\]
The matrix, the order $120$, and the element orders are also
checked independently in
\texttt{code/CD\_E\_identity\_matrix\_counterexample.py}.
Thus its presentation complex is integrally acyclic. It cannot
be aspherical: a group with a finite-dimensional classifying
space is torsion-free, whereas $P$ is finite and nontrivial.
One way to check the latter obstruction directly is to restrict
the putative finite free resolution to a subgroup of prime
order; its periodic integral cohomology is nonzero in
arbitrarily large degrees. Therefore even a successful
nonsingular or integral-unimodular conversion of E.12 would
require a separate geometric proof. The special finite-prefix
elimination in E.7 is precisely what supplies such a proof
for the tree systems in INDLOCAL A.27--A.31.

\paragraph{Statement E.14 (one-pair extensions cannot produce a new simple group).}
Let $X_0\subseteq X_1\subseteq\cdots\subseteq X_\alpha
\subseteq\cdots$ be a continuous well-ordered sequence of
aspherical two-dimensional CW complexes. Suppose each successor
adjoins one one-cell and one two-cell whose attaching word has
exponent sum $\pm1$ in that new one-cell. If the union has
nonabelian simple fundamental group, then every inclusion induces
an isomorphism on fundamental groups.

To see this, every successor is an integral homology equivalence,
as is the inclusion of any stage into the union. Its relative
exponent rows are independent. For any normal subgroup $N$ of a
stage group $H$, the nonsingular theorem applied after replacing
coefficients by $H/N$ gives
\[
 H\cap\langle\!\langle N\rangle\!\rangle_G=N.
\]
Thus every nontrivial stage group is simple when $G$ is simple.
Klyachko's Main Theorem \cite{EKlyachkoSimple2005} states that
for torsion-free $H$ the group $\langle H,t\mid w\rangle$ is
simple precisely when $H$ is simple and $w$ is conjugate in
$H*\langle t\rangle$ to $t^{\pm1}h$ for some $h\in H$.
In that case the inclusion of $H$ is an isomorphism.
If a stage group is trivial, its next one-pair unimodular
extension is trivial as well. Continuity at limits now proves
the assertion by transfinite induction.

This result explains another limit of sequential equation
closure as a proposed construction of a simple group with small
second homology. It does not apply to a simultaneous countable
block whose attaching words depend on arbitrarily later
variables. Such blocks, rather than a well-ordering into
one-pair extensions, are required in E.2 and E.7.

\paragraph{Statement E.15 (small retracts for sparse power-conjugation constructions).}
Let
\[
 G=\langle x_s\ (s\in S)\mid r_s\ (s\in S),\ q\ (q\in Q)\rangle,
\]
where each relator is a finite word and deleting all occurrences
of $x_s^{\pm1}$ from $r_s$ leaves the identity in the free group
on the other generators. If $|Q|\leq\nu$ for an infinite cardinal
$\nu$, then every subset of $G$ of size at most $\nu$ is contained
in a retract of $G$ of size at most $\nu$. In particular, when
$Q$ is countable, every nonidentity element is detected by a
countable quotient.

\emph{Proof.}
Choose words for the specified elements. Start with their
generator supports and the supports of all relators in $Q$.
Close this set under the operation of adjoining the finite
support of $r_s$ whenever $s$ is already present. The closure
$C\subseteq S$ has size at most $\nu$ after countably many
passes. Form the presented group
\[
 G_C=\langle x_s\ (s\in C)\mid r_s\ (s\in C),\ q\ (q\in Q)\rangle.
\]
Its presentation includes only generators from $C$, so there is
a homomorphism $j:G_C\to G$ taking each generator to its namesake.
Conversely, send $x_s$ to its namesake when $s\in C$, and to $1$
otherwise. Every $q$ and every $r_s$ with $s\in C$ is respected.
When $s\notin C$, the stipulated deletion property makes $r_s$
an identity after this substitution, irrespective of the values
of its other letters. Thus the substitution defines
$p:G\to G_C$ with $pj=\mathrm{id}_{G_C}$.
Consequently $j$ is injective, its image contains the specified
elements, and $jp$ is the desired retraction.\hfill$\square$

Here is a useful homological specialization. Suppose a presentation
of $G$ has generators $x_v$ and only relators
\[
 x_u^{-1}x_vx_u=x_v^2,
\]
with any small set of ordered pairs $(u,v)$, allowing repeated
pairs. If its presentation complex is an actual two-dimensional
$K(G,1)$, put
\[
 \nu=\max\{\aleph_0,\operatorname{rank}H_1(G;\mathbb Z),
                         \operatorname{rank}H_2(G;\mathbb Z)\}.
\]
Then every nonidentity element of $G$ has a quotient of size at
most $\nu$, with no exponent $\aleph_0$ on this bound.
Indeed the boundary of a two-cell with target $v$ is $-x_v$.
Choose one relator for every target that occurs. The remaining
relators form a set of cardinality at most $\nu$: subtracting
the chosen relator with the same target gives a basis of the
kernel of the cellular boundary, whose rank is that of $H_2$.
The generators that are never targets form a basis of $H_1$;
place all of them in every closure $C$ above. Each chosen
power-conjugation relator has the deletion property for its
target generator. The same retraction proof therefore applies.
In particular, such an uncountable Higman-type network with
countably generated first and second homology cannot answer E2.
The arbitrary coefficient shifts in A.27 do not have the
deletion property, so this argument does not extend to those
systems.

\paragraph{Statement E.16 (finite homology rank in a perfect normal subgroup).}
Let $G$ admit an actual CW classifying space of dimension at most two,
and let $N\triangleleft G$ be perfect. Suppose that
\[
 H_2(N;\mathbb Z)\cong\mathbb Z^r,
 \qquad 0<r<\infty.
\]
Then the quotient $Q=G/N$ acts faithfully on $H_2(N;\mathbb Z)$,
is finite, and satisfies
\[
                         |Q|\mid r+1.
\]
In particular $r=1$ implies $[G:N]\leq2$, and a perfect $G$ then
equals $N$. For every finite positive $r$, if $G$ has no nontrivial
finite image, then $G=N$. At the distinct boundary $r=0$, the
conclusion is instead $\operatorname{cd}_{\mathbb Z}Q\leq2$;
in particular $Q$ is torsion-free. No geometric-dimension assertion
for this quotient is made.

\emph{Proof.}
Let $X$ be the given two-dimensional classifying space and let
$X_N\to X$ be its regular $N$-cover. Its cellular chains are free
$\mathbb ZQ$-modules, without a finite-rank assumption. Since the
cover is a classifying space for $N$, perfection and the absence
of three-cells give the exact sequence
\[
 0\longrightarrow M\longrightarrow C_2(X_N;\mathbb Z)
 \longrightarrow C_1(X_N;\mathbb Z)
 \longrightarrow C_0(X_N;\mathbb Z)
 \longrightarrow\mathbb Z\longrightarrow0,
 \qquad M=H_2(N;\mathbb Z).
 \tag{E.16.1}
\]
This also explains why $H_2(N;\mathbb Z)$ is free abelian.

First $Q$ is torsion. If $q\in Q$ had infinite order, restrict
the rational cellular complex to
$R=\mathbb Q[\langle q\rangle]=\mathbb Q[t,t^{-1}]$.
Its chain modules remain free: restriction from $\mathbb Q Q$
uses coset representatives as a basis. Because there are no
three-cells, $H_2(X_N;\mathbb Q)$ is a submodule of the free
$R$-module $C_2(X_N;\mathbb Q)$. The ring $R$ is a principal
ideal domain, so this submodule is free, including when the
ambient module has infinite rank. A nonzero free $R$-module
has infinite rational dimension, contradicting
$H_2(X_N;\mathbb Q)\cong\mathbb Q^r$.

Next the action of $Q$ on $M$ is faithful. Suppose its kernel
contains an element of prime order $p$, and let $P$ be the
cyclic subgroup it generates. Restrict (E.16.1) to $\mathbb ZP$.
The three chain modules remain free and have zero positive
$P$-cohomology. This is true for arbitrary free ranks: the
cyclic periodic resolution has finite free terms and its
cochain operators $g-1$ and $1+g+\cdots+g^{p-1}$ are exact
in positive degrees on every direct sum of regular modules.
Three applications of the cohomology long exact sequence give
\[
 H^4(P;M)\cong H^1(P;\mathbb Z)=0.
\]
On the other hand the assumed trivial action on $M=\mathbb Z^r$
gives $H^4(P;M)\cong(\mathbb Z/p)^r\ne0$, by the same periodic
resolution. This contradiction excludes every prime-order
element in the kernel. Since $Q$ is torsion, any nonidentity
kernel element would have a prime-order power. Thus
\[
                       Q\hookrightarrow\operatorname{GL}_r(\mathbb Z).
\]

For completeness, a torsion subgroup of
$\operatorname{GL}_r(\mathbb Z)$ is finite. Reduction modulo
$3$ is injective on such a subgroup because its kernel has
no torsion. To check this last assertion, replace a hypothetical
nonidentity torsion matrix in that kernel by a power of prime
order $\ell$ and write it as $A=I+3^sB$, with $s\geq1$ and
$B$ nonzero modulo $3$. If $\ell\ne3$, expansion of $A^\ell$
modulo $3^{s+1}$ contradicts $A^\ell=I$. If $\ell=3$, expansion
modulo $3^{s+2}$ gives the same contradiction: the leading
term is $3^{s+1}B$, while both later terms vanish modulo
$3^{s+2}$. Hence $Q$ injects into the finite group
$\operatorname{GL}_r(\mathbb F_3)$.

It remains to prove the sharper order divisibility. Let $P$ be
a Sylow $p$-subgroup of the now finite group $Q$, put
$k=\mathbb F_p$ and $A=kP$, and take cellular chains over $k$.
The integral universal coefficient sequence gives
\[
 H_1(X_N;k)=0,
 \qquad H_2(X_N;k)=M\otimes k\cong k^r:
\]
the possible degree-two Tor term is
$\operatorname{Tor}_1^{\mathbb Z}(H_1(X_N;\mathbb Z),k)=0$,
and the degree-one term uses the torsion-free group $H_0=\mathbb Z$.
Consequently the restricted $A$-complex gives an exact sequence
\[
 0\longrightarrow M_k\longrightarrow C_2(X_N;k)
 \longrightarrow C_1(X_N;k)\longrightarrow C_0(X_N;k)
 \longrightarrow k\longrightarrow0,
 \qquad \dim_k M_k=r.
 \tag{E.16.2}
\]

Choose a partial resolution of $k$ by finite free $A$-modules
$D_0,D_1,D_2$, and write $L$ for its third syzygy. These
modules and $L$ are finite-dimensional because $A$ is a
finite-dimensional algebra. The iterated Schanuel lemma
applied to (E.16.2) gives
\[
 M_k\oplus D_2\oplus C_1\oplus D_0
 \cong L\oplus C_2\oplus D_1\oplus C_0.
 \tag{E.16.3}
\]
Here the $C_i$ may still have infinite free ranks. The following
finite stabilization argument prevents cancellation of infinite
dimensions from entering the proof.

If finite $A$-modules $U,V$ become isomorphic after addition
of arbitrary free modules, the inclusion and projection maps
of that isomorphism give maps $f:U\to V$, $g:V\to U$ such
that $gf-1_U$ and $fg-1_V$ factor through free modules.
Because $U,V$ are finitely generated, both factorizations
use finite free summands. Write $gf-1_U=ba$ through a finite
free module $F$. Then
\[
 j=(f,a):U\longrightarrow V\oplus F,
 \qquad s=(g,-b):V\oplus F\longrightarrow U
\]
satisfy $sj=1_U$. The complement $K=\ker s$ is finite.
Moreover $1-js$ factors through a finite projective module:
its upper-left block is $1-fg$, and each other block either
has domain or range $F$. Restricting this factorization to
$K$ factors $1_K$ through a finite projective module, so
$K$ is itself projective. We have therefore obtained an
actual finite stabilization
\[
                         V\oplus F\cong U\oplus K.
\]
For a finite $p$-group, $A=kP$ is local and every finite
projective $A$-module is free. One can see locality by
inducting on $|P|$: a central element of order $p$ makes
the ideal it generates minus $1$ nilpotent, and the quotient
is the group algebra of the smaller $p$-group; thus the
augmentation ideal is nilpotent and $A$ has residue field $k$.
Lifting a basis modulo that ideal and applying Nakayama's
lemma proves the asserted freeness of finite projectives.
Thus the finite stabilization shows
\[
                         \dim_k U\equiv\dim_k V\pmod{|P|}.
\]

Apply this to (E.16.3). The finite partial resolution gives
\[
 \dim_k L=\dim_k D_2-\dim_k D_1+\dim_k D_0-1
          \equiv-1\pmod{|P|}.
\]
Hence $r\equiv-1\pmod{|P|}$. Each Sylow subgroup order
divides $r+1$, and their orders are pairwise coprime, proving
$|Q|\mid r+1$.

When $r=1$, this gives $Q=1$ or $Q\cong C_2$; a quotient
of a perfect group is perfect, so a perfect $G$ allows only
the first case. The assertion when $G$ has no nontrivial
finite image follows directly from finiteness of $Q$.
Finally, if $r=0$, (E.16.1) itself is a free
$\mathbb ZQ$-resolution of $\mathbb Z$ of length at most two.
It proves $\operatorname{cd}_{\mathbb Z}Q\leq2$, and the
cyclic periodic resolution excludes nontrivial finite-order
elements in $Q$.\hfill$\square$

This restriction does not settle E2. A countable homology
core need not be normal, and its normal closure need not
retain the finite second-homology rank. In particular, adding
unrestricted auxiliary conjugators and then taking a normal
closure still requires an independent homology calculation.


\paragraph{Corollary E.16a (the cost of taking a normal closure).}
Let $Y\subseteq X$ be an actual subcomplex of a two-dimensional
classifying space, with injective $H=\pi_1Y\leq G=\pi_1X$,
and suppose $Y$ is aspherical, $H$ is perfect, and
$H_2(Y;\mathbb Z)\cong\mathbb Z^{(\kappa)}$.
Put $N=\langle\!\langle H\rangle\!\rangle_G$ and $Q=G/N$.
Then $N$ is perfect and there is an injective $\mathbb ZQ$-module map
\[
                   (\mathbb ZQ)^{(\kappa)}\hookrightarrow H_2(N;\mathbb Z).
\]
In particular $\operatorname{rank}H_2(N)\geq\kappa\,|Q|$.

\emph{Proof.}
The normal closure is generated by conjugates of the perfect
subgroup $H$, so its abelianization is zero. In the regular
$N$-cover $X_N\to X$, the inverse image of $Y$ is the disjoint
union of copies $Y_q$ indexed by $q\in Q$: the restricted
covering has trivial monodromy because $H\leq N$. Their cellular
two-cycles inject into those of $X_N$. There are no three-cells
in either space, so the same inclusion is injective on second
homology. Deck transformations permute the copies regularly,
giving exactly the displayed $\mathbb ZQ$-module injection.
Finally $X_N$ is a classifying space for $N$.\hfill$\square$

The conclusion uses an actual subcomplex, not just an abstract
subgroup or a map inducing an isomorphism on ordinary homology.
It quantifies the homology increase that can occur when auxiliary
free conjugators are added and the original perfect group is
then replaced by its normal closure.

\paragraph{Statement E.17 (a fixed-ball extension of the countable diagonal).}
Let $E$ be a countably infinite nonabelian simple group, and suppose
\[
 N(g):=\sup_{a\ne1}\rho_a(g)<\infty\qquad(g\in E),
\]
where $\rho_a$ is conjugacy length with respect to $a^{\pm1}$.
Assume there are $d\ne1$ and an integer $L\geq0$ such that, writing
$D=N(d)$, every conjugate of every $c\ne1$ with $N(c)\leq2D$
can be reached by a conjugator $v$ satisfying $N(v)\leq L$.
Then the conclusion of the killing-block lemma A.27 holds for $E$:
for any actual aspherical CW complex $B$ of dimension at most two,
$K\leq H=\pi_1B$, and epimorphism $q:K\twoheadrightarrow E$, there
is an actual aspherical extension $B\subseteq Y$ of dimension at
most two such that
\[
 H\hookrightarrow P=\pi_1Y,\qquad
 H_*(B;\mathbb Z)\cong H_*(Y;\mathbb Z),\qquad
 \langle\!\langle H\rangle\!\rangle_P=P,
\]
and no homomorphism $P\to E$ restricts to $q$ on $K$.
Only countably many cells are attached.

The hypothesis is weaker than requiring one transport bound on
all of $E$: it concerns just the ball of radius $2N(d)$.
In particular it follows from the bounded-conjugator condition
stated as Assumption~4.1 in
\cite{EGismatullinMajcherZiegler2024}, applied to the constant
family $(E,N)$. No ultraproduct theorem from that paper is used.

\emph{Proof.}
The function $N$ is integer-valued, symmetric, conjugation-invariant,
and subadditive: each of these assertions follows from the
corresponding property of every $\rho_a$ by taking the supremum.
Put $C=L+D$, $m=2C$, and
$B_C=\{g\in E:N(g)\leq C\}$. This ball contains $1$ and $d$.
For $a\in E$ define
\[
 W_a(u_1,v_1,\ldots,u_m,v_m)
       =\prod_{i=1}^{m}v_i[a,u_i]v_i^{-1}.
 \tag{E.17.1}
\]
We claim that for every $a\ne1$,
\[
             B_{2C}\subseteq W_a(B_C^{\,2m}).
 \tag{E.17.2}
\]
Choose a conjugate $u$ of $d$ with $c=[a,u]\ne1$.
Such a conjugate exists because the conjugates of $d$ normally
generate the simple group $E$; if all commuted with $a$, then
$a$ would be central. We have $N(u)=D$ and
$N(c)\leq2D$. Each $g\in B_{2C}$ is a product of at most
$\rho_c(g)\leq N(g)\leq2C=m$ conjugates of $c^{\pm1}$.
For every factor, the assumed transport bound replaces its
conjugator by one with $N$ at most $L$. Positive factors now
have the form in (E.17.1), with $N(u_i)=D$ and $N(v_i)\leq L$.
For a negative factor use
\[
                  c^{-1}=u[a,u^{-1}]u^{-1}.
\]
Its new conjugator has norm at most $L+D=C$, and its commutator
variable still has norm $D$. Set unused variables equal to $1$.
All variables lie in $B_C$, proving (E.17.2). Consequently,
for every $a\ne1$ and $b\in B_C$,
\[
                  B_C\subseteq W_a(B_C^{\,2m})b,
 \tag{E.17.3}
\]
since $N(gb^{-1})\leq2C$ whenever $g,b\in B_C$.
Notice that no bound on $N(a)$ is required.

Enumerate $E$ without repetition as $e_1,e_2,\ldots$ and
choose $h\ne1$ in $E$. Use one variable $x_s$ at each vertex
of the rooted $2m$-ary tree, with root $o$. We shall choose
coefficients $b_s\in B_C$, starting with $b_o=1$, and impose
\[
 \begin{aligned}
 x_o&=W_h(x_1,\ldots,x_{2m})b_o,\\
 x_s&=W_{x_oe_n^{-1}}(x_{s1},\ldots,x_{s,2m})b_s
            &&(|s|=n\geq1).
 \end{aligned}
 \tag{E.17.4}
\]
After coefficients at depths less than $n$ have been fixed,
substitute $x_o=e_n$ into those earlier right sides. The
nonroot coefficients $e_ne_j^{-1}$, $j<n$, are nonidentity.
The resulting map $T_n$ from depth-$n$ values to the computed
root value satisfies
\[
                  B_C\subseteq T_n(B_C^{(2m)^n}).
 \tag{E.17.5}
\]
Indeed, lift any prescribed root value in $B_C$ backwards
through the finite tree, using (E.17.3) at each vertex. The
lists of child variables are disjoint, so the choices do not
interfere and all chosen intermediate values remain in $B_C$.
As $B_C$ has at least two elements, choose a value in $B_C$
different from $e_n$ and take a corresponding tuple
$(b_s)_{|s|=n}\in B_C^{(2m)^n}$ in (E.17.5).

A solution in $E$ would have $x_o=e_n$ for some $n$.
At depth $n$ every $W$ has coefficient $1$, so (E.17.4)
forces $x_s=b_s$ there. Earlier equations then force
$x_o=T_n((b_s))\ne e_n$, a contradiction. Thus this countable
system has no solution in $E$.

Lift every chosen coefficient through $q$ into $K\leq H$
and attach the displayed variable loops and relator two-cells
to $B$. Each $W$ has exponent zero in its coefficient and
in every child variable, so the relative cellular boundary
sends each relator precisely to its own variable. It is the
identity on the countable free abelian direct sums, proving
the integral homology isomorphism.

For the subcomplex containing the cells down to depth $n$,
eliminate the nonroot generator--relator pairs in descending
order of depth. Only free frontier loops, the root loop,
and one root relator remain over $B$. That last relator has
root exponent one, since the outer word $W_h$ has exponent
zero in every child. The coefficient group is the torsion-free
free product of $H$ with the remaining free group.
Forester--Rourke's unimodular adjunction theorem, in the
form audited in E.7, gives an aspherical finite-depth
subcomplex and an injective coefficient map. This invocation
uses \cite[Main Theorem]{EForesterRourkeDiagrams}; it does not
assume that an arbitrary identity-exponent system is aspherical.
Every sphere or nullhomotopy in the increasing CW union is
contained in a finite-depth subcomplex, so $Y$ is aspherical
and $H\to\pi_1Y$ is injective.

Finally, killing $H$ makes the root coefficient and $b_o$
trivial, so $x_o=1$. All other coefficients then become
trivial, and (E.17.4) makes every variable $1$. Hence $H$
normally generates $P$. A homomorphism $P\to E$ extending
$q$ would supply the impossible solution above.
\hfill$\square$

The same enumeration and $\omega_1$-stage bookkeeping as in
A.28 can therefore exclude these additional targets while
preserving the actual two-dimensional model and its homology.
This does not settle E2. A remaining countable simple target
with pointwise finite $N$ must now fail the stated transport
bound for every $d\ne1$: for every $L$ there must be some
$c\ne1$ with $N(c)\leq2N(d)$ and a conjugate of $c$ that
has no conjugator of norm at most $L$. The elements $c$ may
vary with $L$; no fixed-conjugacy-class assertion is inferred.

\paragraph{Statement E.18 (normal subgroup homology bounds the index).}
Let $X$ be an actual two-dimensional CW classifying space for a
group $G$, let $N\trianglelefteq G$, and suppose
\[
       0\ne H_2(N;\mathbb Z)\cong\mathbb Z^{(\kappa)}.
\]
If $\kappa$ is finite, then $[G:N]$ is finite.
If $\kappa$ is infinite, then $[G:N]\leq\kappa$.
No perfectness hypothesis is required.

\emph{Proof.}
Put $Q=G/N$ and let $X_N\to X$ be the corresponding regular
cover. Because there are no three-cells,
\[
 M:=H_2(N;\mathbb Z)=H_2(X_N;\mathbb Z)
       =\ker\bigl(C_2(X_N)\to C_1(X_N)\bigr)
\]
is literally a $\mathbb ZQ$-submodule of $C_2(X_N)$.
Choose lifts and orientations of the two-cells of $X$.
They give a basis $\mathcal E$ of $C_2(X_N)$ that is a
disjoint union of free $Q$-orbits. Choose a free abelian
basis $(m_i)_{i\in I}$ of $M$, with $|I|=\kappa$, and set
\[
                  S=\bigcup_{i\in I}\operatorname{supp}_{\mathcal E}(m_i).
\]
The set $S$ is nonempty. It is finite when $I$ is finite,
and has cardinality at most $\kappa$ when $\kappa$ is
infinite. Moreover it is $Q$-invariant. Indeed, for every
$q\in Q$ each $qm_i$ is a finite integral combination of
the $m_j$, so $q\operatorname{supp}(m_i)\subseteq S$.
Taking unions gives $qS\subseteq S$, and applying this to
$q^{-1}$ gives equality. Since every $Q$-orbit in
$\mathcal E$ is free, any one orbit in $S$ has cardinality
$|Q|$. Thus $|Q|\leq|S|$, which gives both conclusions.
\hfill$\square$

\paragraph{Corollary E.18a (faithfulness for arbitrary positive rank).}
Under the preceding hypotheses, if $N$ is perfect, the
action $Q\to\operatorname{Aut}(M)$ is faithful, for
arbitrary nonzero rank $\kappa$.

\emph{Proof.}
For any nonzero $m\in M$, its stabilizer in $Q$ permutes
the finite nonempty support of $m$ in the free $Q$-set
$\mathcal E$, and is therefore finite. In particular the
kernel $K$ of the action on $M$ is finite. If $K\ne1$,
Cauchy's theorem supplies a cyclic subgroup $C_p\leq K$.
Since $N$ is perfect, the cellular sequence
\[
 0\longrightarrow M\longrightarrow C_2(X_N)
  \longrightarrow C_1(X_N)\longrightarrow C_0(X_N)
  \longrightarrow\mathbb Z\longrightarrow0
\]
is exact. Restriction to $C_p$ leaves all the $C_j$ free,
of possibly infinite rank. The usual two-periodic
resolution of $C_p$, whose terms are finitely generated,
shows that these free modules have zero positive
cohomology even at infinite rank. Dimension shifting gives
\[
 H^4(C_p;M)\cong H^1(C_p;\mathbb Z)=0.
\]
But the action on the nonzero free abelian group $M$ is
trivial, so $H^4(C_p;M)=M/pM\ne0$, a contradiction.
\hfill$\square$

\paragraph{Statement E.19 (homology of the literal Corson--Obraztsov construction).}
Let $\kappa$ be an uncountable cardinal admitting a J\'onsson algebra,
and choose the group $J_\kappa$ by Corson's proof
\cite[Proposition~5 and proof of Theorem~1]{ECorson2022Jonsson},
using the presentation in Obraztsov's Theorem~A
\cite[pp.~268--269, 276--278]{EObraztsov1996Scheme} with $H=1$, $m=\infty$, and infinite cyclic
factors indexed by $\kappa$. Then
\[
                 H_2(J_\kappa;\mathbb Z)\cong\mathbb Z^{(\kappa)}.
\]
In particular this specific J\'onsson construction does not give an
acyclic group or a group with countable second homology.

\emph{Proof.}
Write $F=*_{\alpha<\kappa}\langle x_\alpha\rangle$ and
$J_\kappa=F/N$, where $N$ is normally generated by the defining
relators in Obraztsov's construction. Let $\mathcal R$ contain one
oriented representative from each free-product conjugacy class up to inversion
of defining relators (equivalently, cyclic classes after cyclic reduction);
retain the displayed representatives used below. His notion of
diagrammatic asphericity, stated on p.~268 and established on p.~285,
has the following sufficient consequence:
\begin{equation}\label{E:Jonsson:relationcoinvariants}
                  N/[F,N]\cong\mathbb Z^{(\mathcal R)}.
\end{equation}
The convention is the one made explicit by Muranov
\cite[Section~6, Definition~60; Sections~3.3--3.4]{EMuranov2007Width}: a spherical diagram splits into elementary two-face
spheres by diamond moves; his Definitions~33 and 38--40 describe
these moves and cancellations. Section~6 expressly identifies this
with Ol'shanskii's diagrammatic asphericity. In an oriented sphere,
diamond moves retain the oriented face boundary labels, and the
two faces of an elementary sphere have inverse boundary labels.
To see the consequence directly,
the displayed map is onto by normal generation. A relation in its
kernel can, after writing the relevant elements of $[F,N]$ as
commutators and expanding their $N$-entries into conjugates of
defining relators, be written as an identity product of conjugates
of relators. Its signed multiplicities are the originally proposed
relation. Draw the corresponding spherical identity diagram.
Removal of a cancelling pair of defining cells changes none of
these signed multiplicities. Repeated diagrammatic reductions end
with factor cells only; those have no defining-relator multiplicities.
Thus all the originally proposed multiplicities are zero. The
same argument works after expanding each letter of a cyclic factor
as a power of $x_\alpha$; factor cells fill in its circle and contribute
zero multiplicity. Every diagram involved is finite, so no
cardinality restriction enters this argument.

The five-term sequence for $1\to N\to F\to J_\kappa\to1$,
or equivalently Hopf's formula, now identifies
\begin{equation}\label{E:Jonsson:Hopfkernel}
 H_2(J_\kappa;\mathbb Z)
 =\ker\!\left(\mathbb Z^{(\mathcal R)}
                 \xrightarrow{\partial}\mathbb Z^{(\kappa)}\right),
\end{equation}
where $\partial$ records ordinary exponent sums.

For each pair of distinct cyclic factors with indices $\alpha<\beta$,
the special periods listed on p.~277 include
$A=a[a,c]^t$, with nonidentity $a\in\langle x_\alpha\rangle$ and
$c\in\langle x_\beta\rangle$ and a permitted positive integer $t$.
Choose $a,c$ in the oriented halves used there; if the auxiliary
choice $a<b$ is needed, choose $a$ not maximal in its infinite
oriented half. Relation~(2.2) on that page includes, in particular,
the three relators with $k=6,7,8$ in the family
\[
 R_k=c^{-1}A^n
       \prod_{j=0}^{h-2}\bigl(c A^{n+30j+k}\bigr).
\]
Their exponent sums are
\[
 \partial R_k
 =\bigl(hn+15(h-2)(h-1)+(h-1)k\bigr)[a]+(h-2)[c].
\]
Consequently
\[
                    z_{\alpha\beta}=R_6-2R_7+R_8
\]
is a nonzero element of the kernel in
\eqref{E:Jonsson:Hopfkernel}. The three relators are distinct:
their exponent sums in $x_\alpha$ are distinct, because $a$ is a
nonzero power of that generator. Their positive $x_\alpha$-coefficient
relative to $[a]$, together with the fixed nonzero $[c]$ coefficient,
also distinguishes their chosen orientations from inverse relators.

Partition the factor index set into $\kappa$ disjoint pairs.
The corresponding $z_{\alpha\beta}$ have disjoint supports in the
free abelian group on the relators, and are therefore independent.
Indeed every defining word in such a support uses exactly the
two factors of that pair. This proves that the rank of $H_2$ is
at least $\kappa$. There are at most $\kappa$ finite words over
the alphabet of the factors, hence at most $\kappa$ defining
relators. Formula~\eqref{E:Jonsson:Hopfkernel} makes $H_2$ a
free abelian group of rank at most $\kappa$. Equality follows.
\hfill$\square$

\emph{Scope.}
This calculation uses the explicit presentation selected in
Corson's proof; it does not assert that every J\'onsson group has
large second homology. Nor does it identify Obraztsov's terminology
``diagrammatically aspherical'' with the existence of an actual
two-dimensional classifying space: his terminology also applies
to finite-exponent variants. The signed-relator argument and
Hopf's formula give the claimed obstruction without requiring
that additional geometric assertion.

The finite word calculation was independently checked in
\texttt{code/CD\_E\_Jonsson\_cycles.py} on 108 parameter choices,
including free reduction and the second-difference exponent identity.
The actual proof above uses the symbolic formula with the source's
admissible large parameters, not those small test values.


\paragraph{Statement E.20 (a piecewise linear Higman exclusion).}
Let $\operatorname{BPL}_+(\mathbb R)$ be the group of
orientation-preserving piecewise linear homeomorphisms of the line
whose breakpoints are locally finite and which are bi-Lipschitz.
The bi-Lipschitz constant may depend on the element.
Every homomorphism from the Higman group $H$ to this group is
trivial. Consequently, every group normally generated by images
of copies of $H$, including the group in A.28, has no nontrivial homomorphism
to any subgroup of $\operatorname{BPL}_+(\mathbb R)$.

\emph{Proof.}
First suppose that $uvu^{-1}=v^2$ in this group and that $v\ne1$
has a fixed point. A nonempty component of its support has a finite
endpoint $p$. On the side pointing into that component, local
finiteness of the breakpoints makes $v$ linear near $p$, with
one-sided slope $\lambda\ne1$. Indeed, $v(p)=p$, so slope one
would make that interval pointwise fixed. The map $u$ preserves
$\operatorname{Fix}(v)$ and its complementary components, because
$\operatorname{Fix}(v^2)=\operatorname{Fix}(v)$.
At every fixed point, the one-sided derivatives of $u$ and $u^{-1}$
cancel in differentiating $uvu^{-1}$. Orientation preserves the
chosen side. The relation therefore gives
\[
 v'_{\pm}(up)^2=v'_{\pm}(p),\qquad
 v'_{\pm}(u^{-k}p)=\lambda^{2^k}\quad(k\ge0).
\]
For $\lambda>1$ these slopes are unbounded, and for $0<\lambda<1$
they tend to zero. Both alternatives contradict the finite upper
and positive lower slope bounds supplied by the bi-Lipschitz
condition on the single element $v$. Thus a nontrivial $v$ in such
a relation has no fixed point.

Now suppose that
\[
 aba^{-1}=b^2,\qquad cac^{-1}=a^2.
\]
If $b\ne1$, it is fixed-point-free. An increasing topological change
of coordinates takes $b$ to $t\mapsto t+1$ or $t\mapsto t-1$:
choose a homeomorphism on a fundamental interval and extend using
the iterates of $b$. In either case the first relation gives
$a(t+1)=a(t)+2$ in these coordinates. Hence
\[
 a(t+n)-(t+n)=a(t)-t+n\qquad(n\in\mathbb Z).
\]
Continuity supplies a fixed point of $a$. Fixed points do not depend
on the chosen coordinates. Applying the preceding argument to the
second relation in the original coordinates gives $a=1$, and then
$b=b^2$ gives $b=1$, a contradiction. Every Higman generator is the
$b$-term of two consecutive such relations, after replacing all
named generators by their inverses if using the presentation in
E.4. Thus every generator image is trivial. Normal generation gives
the final assertion. \hfill$\square$

The local finiteness and global bi-Lipschitz hypotheses are part
of the assertion. In particular, this is not an exclusion of
arbitrary groups of orientation-preserving homeomorphisms of the
line. A finite global set of positive slopes is sufficient for the
bi-Lipschitz condition. The proof uses two consecutive Higman
relations; a single power-conjugacy relation can hold nontrivially
for affine maps of the line.

\paragraph{Statement E.21 (a countable-quotient block with one homology cost).}
Let $P$ be a perfect group with an actual two-dimensional
classifying space $X$, and let $q:P\twoheadrightarrow E$ be a
surjection onto a countable group. Choose $h\in P$ with $q(h)\ne1$.
There is an injective homomorphism $P\hookrightarrow G$, realized
by a subcomplex inclusion between actual two-dimensional
classifying spaces, such that
\[
 H_1(G;\mathbb Z)=0,\qquad
 H_2(G;\mathbb Z)\cong H_2(P;\mathbb Z)\oplus\mathbb Z,
 \tag{E.21.1}
\]
and no homomorphism $G\to E$ extends $q$. The construction adds
only countably many cells to $X$, and $G$ is normally generated
by the image of $P$.

\emph{Proof.}
Choose $c_0,c_1,\ldots\in P$ whose $q$-images enumerate $E$;
repetitions are allowed. Use the convention $[x,y]=xyx^{-1}y^{-1}$.
Start from $P*\langle a\rangle$ and set
\[
 L=\left\langle P,a,y_0,y_1,\ldots\ \middle|\
 h=[a,y_0],\quad
 y_n=[a c_n^{-1},y_{n+1}]\ (n\ge0)\right\rangle.
 \tag{E.21.2}
\]
This presentation is an increasing union of injective cyclic HNN
extensions, as follows. The first new relation is
\[
 y_0 a^{-1}y_0^{-1}=a^{-1}h.
\]
Both associated cyclic subgroups are infinite: each displayed
endpoint has nonzero $a$-exponent in $P*\langle a\rangle$.
After $y_n$ has been added, write $k_n=a c_n^{-1}$. The next
relation is
\[
 y_{n+1}k_n^{-1}y_{n+1}^{-1}=k_n^{-1}y_n.
\]
The left endpoint is nontrivial of infinite order in the embedded
initial free product. The right endpoint has exponent one in the
most recently added stable letter $y_n$. The homomorphism from
that finite HNN stage to $\mathbb Z$ which sends $y_n$ to one and
its base group to zero proves that this endpoint also has infinite
order. The HNN normal form theorem therefore embeds each stage
in the next one.

Use $X\vee S^1$ for the initial vertex space and attach a cylinder
on a circle for each cyclic HNN extension. All vertex spaces are
aspherical and all edge maps induce injections, so the usual tree
of universal covers proves asphericity at each finite stage.
These are actual two-dimensional CW spaces, with the previous
stage as a subcomplex. A sphere in the union has image in a finite
subcomplex and hence in some finite stage; it is null-homotopic
there. The union is therefore an actual two-dimensional $K(L,1)$
containing $X$.

For ordinary relative cellular chains over $X$, there is one new
one-cell $a$ and the one-cells $y_n$. The root two-cell
$h=[a,y_0]$ has zero relative boundary. Each subsequent two-cell
$y_n=[a c_n^{-1},y_{n+1}]$ has relative boundary $y_n$, up to the
choice of orientation: all occurrences of $a$ and $y_{n+1}$ have
cancelling exponent sums. Thus
\[
 H_2(K(L,1),X;\mathbb Z)=\mathbb Z,\qquad
 H_1(K(L,1),X;\mathbb Z)=\mathbb Z\langle a\rangle.
\]
The connecting homomorphism from this relative $H_2$ sends its
generator to the class of $h$ in $H_1(P;\mathbb Z)$, which is
zero by perfection. Since there are no relative three-cells, the
long exact sequence gives
\[
 H_2(L;\mathbb Z)\cong H_2(P;\mathbb Z)\oplus\mathbb Z,
 \qquad H_1(L;\mathbb Z)=\mathbb Z\langle a\rangle.
\]
The splitting is not claimed to be canonical.

To kill the remaining first homology while retaining normal
generation, first adjoin a stable letter $t_0$ with
\[
 t_0 a t_0^{-1}=h.
\]
Both endpoints have infinite order: $a$ lies in the embedded
initial free product, and $h\ne1$ lies in the torsion-free group
$P$. Next adjoin $t_1,t_2,\ldots$ with
\[
 t_n=[h,t_{n+1}]\qquad(n\ge0).
\]
These again are injective cyclic HNN extensions, since the $n$th
relation is
\[
 t_{n+1}h^{-1}t_{n+1}^{-1}=h^{-1}t_n.
\]
The left endpoint has infinite order in the embedded $P$; the right
endpoint has exponent one in the latest stable letter $t_n$.
Let $G$ be the union. The same graph-of-spaces argument supplies
an actual two-dimensional classifying space including $X$.
The first new HNN edge kills the primitive first-homology class
$a$ and introduces $t_0$. Each later edge kills $t_n$ and introduces
$t_{n+1}$. The edge map on first homology is injective at each of
these steps, so Mayer--Vietoris leaves $H_2$ unchanged. Taking the
union kills every such first-homology class. This proves (E.21.1).
Equivalently, in the full ordinary relative cellular chain complex,
the added boundaries are $a$ and then $t_n$; together with the
previous $y_n$ boundaries, they make every new one-cell a boundary
and leave exactly the one root two-cycle.

If $f:G\to E$ extended $q$, choose $n$ with
$f(a)=q(c_n)$. Then $f(k_n)=1$, so the $n$th tail relation gives
$f(y_n)=1$. Successive earlier tail relations imply
$f(y_{n-1})=\cdots=f(y_0)=1$. The root relation now gives
$q(h)=[f(a),f(y_0)]=1$, a contradiction.

Finally, quotient by the normal closure of $P$. Then $h=1$,
so $t_0 a t_0^{-1}=h$ gives $a=1$. The relations for the $y_n$
now give $y_n=1$ for every $n$, and the relations
$t_n=[h,t_{n+1}]$ give $t_n=1$ for every $n$.
Thus the quotient is trivial, proving normal generation.
\hfill$\square$

Each use of this block adds an independent integral second-homology
class. In an increasing construction by these actual subcomplex
inclusions, subsequent two-dimensional attachments cannot remove
any existing cellular two-cycle. Thus the block proves a normally generating quotient
obstruction with an exact cost; it does not give a homology-preserving
replacement for A.27 or settle E2. If instead the root $h$ belongs
to a fresh Higman copy, the weaker conclusion is merely that this
fresh copy must map trivially whenever the image of $a$ lies in the
image of $P$. The old-coefficient HNN construction above is what
makes the prescribed extension $q$ itself impossible.


\paragraph{Statement E.22 (a locally finite tree Higman exclusion).}
Let
\[
 H=\langle a_0,a_1,a_2,a_3\mid
       a_{i+1}a_i a_{i+1}^{-1}=a_i^2\ (i\bmod4)\rangle.
\]
Every homomorphism from $H$ to the automorphism group of a locally
finite simplicial tree is trivial. Consequently, every group normally generated by images
of copies of $H$, including the group in A.28, has no nontrivial homomorphism
to the automorphism group of a locally finite tree.

\emph{Proof.}
Subdivide edges once to remove inversions. A hyperbolic tree isometry
has positive translation length, and its square has twice that length.
Since each $a_i$ is conjugate to its square, every $a_i$ is elliptic.

We first show that each adjacent pair $a=a_i$, $b=a_{i+1}$ has a
common fixed vertex. Put $A=\operatorname{Fix}(a)$, a nonempty convex
subtree. The relation gives
\[
 A\subseteq\operatorname{Fix}(a^2)=bA.
\]
Choose $p\in\operatorname{Fix}(b)$ and let $z$ be the unique nearest
point of $A$ to $p$. The nearest point of $bA$ to $p$ is $bz$, and
$d(p,bz)=d(p,z)$. Since $z\in A\subseteq bA$, uniqueness of the
nearest point in $bA$ gives $bz=z$. Thus $z$ is fixed by both $a$ and $b$.
In the subdivided simplicial tree the point $z$ may be chosen a vertex.

An elliptic automorphism of a locally finite tree has finite orbits
on vertices: a fixed vertex centers finite invariant balls. Suppose
some $a_i$ is nontrivial, and let $p$ be the smallest prime dividing
the length of any nontrivial vertex orbit of any of the four generators.
Choose $i$ and a vertex orbit of $a_i$ whose length is divisible by $p$.
The adjacent pair $a_i,a_{i+1}$ fixes a vertex $z_i$. A sufficiently
large finite ball about $z_i$ contains this orbit and is invariant under
both generators. Let $\bar a,\bar b$ be their permutations of this ball,
and let $m$ be the order of $\bar a$. Then $p\mid m$ and
\[
 \bar b\bar a\bar b^{-1}=\bar a^2.
\]
Consequently $m$ is odd. The order of the automorphism
$\bar a\mapsto\bar a^2$ of the cyclic group of order $m$ divides
the order of $\bar b$. In particular the multiplicative order
$\operatorname{ord}_p(2)$ divides the order of $\bar b$.
Since $p$ is odd, $\operatorname{ord}_p(2)>1$ and it divides $p-1$.
Choose a prime $q$ dividing $\operatorname{ord}_p(2)$; then $q<p$.
The order of a finite permutation is the least common multiple of
its cycle lengths, so $\bar b$ has a vertex orbit whose length is
divisible by $q$. This contradicts the definition of $p$.
Thus all four generators act trivially.
\hfill$\square$

Local finiteness is essential to this proof. It does not claim that
$H$ has Serre's property FA: its familiar amalgam decomposition does
give a nontrivial action on a tree with infinite valence.


\paragraph{Unresolved step and approaches that do not close it.}
Replacing the reduced power in E.2 by a countable
universal target requires solutions to countable
systems, not merely finite ones; A.27 rules out
the natural countable existentially closed
targets. Repeating the metric argument of
A.30 with the topology from A.29 loses the
uniform conjugation bound. Pointwise finite
$N_E$ in A.31 permits uniform bounds on the number of
conjugate factors for reaching any specified finite set,
but gives no bound on the norms of their conjugators.
Statement E.17 closes the tree argument when a transport
bound on one ball is available. Its failure branch remains:
nested tree preimages need not admit a common norm bound.
Finally, replacing old two-cells by the
Bartholdi--Mikhailov blocks can retain
two-dimensional models at each finite step,
but those structural maps are not subcomplex
inclusions: the old cell is removed by a
collapse with a three-cell. A mapping telescope
then has dimension three. These are exact
obstructions to the attempted arguments,
not proofs of nonexistence in E2.
